\chapter{Justificación}

El proyecto se justifica desde una perspectiva técnica, académica y funcional. Desde el punto de vista técnico, la disponibilidad de datasets públicos anotados y de modelos de segmentación médica permite construir prototipos reproducibles sin depender inicialmente de datos privados. SPIDER ofrece una base alineada con el MVP porque incluye RM lumbares sagitales y máscaras de referencia para vértebras, discos intervertebrales y canal espinal; el conjunto de Sudirman / Al-Kafri aporta, de forma complementaria, cortes axiales anotados sobre los niveles discales inferiores; y RSNA / LumbarDISC se suma como dataset complementario empleado en la línea degenerativa por nivel, sin constituir una base de segmentación anatómica primaria. Junto con arquitecturas como U-Net o nnU-Net, este contexto permite enfocar el aporte en la integración de un flujo completo, no solo en el entrenamiento de un modelo aislado.

Desde el punto de vista académico, la propuesta combina procesamiento de imágenes, inteligencia artificial, diseño de software, validación técnica e interacción con usuarios expertos. Esta articulación resulta pertinente para un Proyecto Final de Ingeniería porque exige traducir antecedentes técnicos en un sistema funcional, documentado y evaluable.

Desde el punto de vista funcional, el prototipo puede aportar valor como soporte al análisis estructural de RM lumbar. Las máscaras permitirán visualizar regiones de interés, las mediciones aportarán datos geométricos acotados y la salida editable organizará resultados para revisión profesional. La herramienta se mantiene dentro de un enfoque asistivo: el usuario conserva el control sobre interpretación, corrección y uso final de la información. El estado del arte muestra avances en datasets, modelos, herramientas abiertas y productos comerciales, pero esos elementos suelen aparecer separados o dentro de soluciones cerradas. La brecha que justifica el PFI está en integrar segmentación, mediciones simples, visualización, salida editable y validación técnica/cualitativa en un prototipo académico reproducible.

Finalmente, trabajar con datos públicos reduce riesgos de privacidad durante esta etapa. En una evolución posterior, en la que se contemple conservar estudios de pacientes para habilitar una historia clínica y el seguimiento longitudinal, el tratamiento de datos deberá ajustarse a la normativa vigente (en particular la Ley N.º 25.326 de Protección de los Datos Personales), con almacenamiento seguro, acceso restringido y seudonimización. Para el alcance actual, la justificación se concentra en demostrar factibilidad técnica y utilidad funcional dentro de un MVP responsable.
